\section{Bounded cut cores and wholly prismatic components}
\label{sec:cut-products}

A cut can have exponentially many connected pieces even when its input
coordinates are short. The chamber count alone does not distinguish which
pieces require further geometric work. We now mark a linear number of local
chambers that are not standard normal-disc prisms. Every cut component
avoiding these marks is an interval bundle. This yields a source-certified
count of a bounded conservative core and the remaining wholly prismatic
pieces, including twisted bundles.

\subsection{What the published cutting algorithms already supply}

The earlier wording of Section~\ref{sec:missing} was too broad about the
published record. Lackenby--Yazdi's Theorem 8.6\cite{lackenby-yazdi-genus}
recalls a polynomial algorithm for an orientable normal cutting surface:
it constructs the handle core outside the parallelity bundle, determines
base-surface genus and boundary count, distinguishes product from twisted
bundles, and locates their vertical attachments. The running time is
polynomial in $T\log w(F)$. The source algorithm is Lackenby's
Theorem 9.3\cite{lackenby-parallelity}. It distinguishes the handle core
from the interval-bundle information instead of expanding every handle.

These are published component algorithms, not absent theorems. Their
source-relative realization and integration with the accelerated hierarchy's
full patterns and restart bounds remain missing in our implementation.
Our classification below is smaller than that parallelity-bundle algorithm:
it detects whole components made entirely from prisms, not bundles attached
inside a core component. It does not produce vertical attachment locations
or a core handle structure.


The output contract matters for arbitrary binary multiplicities. A cut
along $m$ parallel meridian discs has $m$ connected pieces. Any ordinary
explicit tetrahedron list for the entire cut therefore contains at least
$m$ tetrahedra. Thus no bound such as $200T$ can describe that explicit
output for all multiplicities on a fixed $T$-tetrahedron source. The printed
sphere/disc retriangulation statements in Theorems 8.2--8.3 require a clarified
restriction or compressed-repetition contract before they can be applied
literally to such inputs. This is a direct output-cardinality observation,
not a criticism of applications to bounded nonparallel cutting families.
We do not use an unrestricted linear explicit-output assertion in the native
algorithm or its complexity proof.

\subsection{A linear set of exceptional chambers}

In the chamber indexing of Section~\ref{sec:cut-complement}, mark the two
endpoints of each occupied quadrilateral chain, coalescing them when $q=0$.
Also mark the innermost chamber of each nonempty triangle arm. There are at
most six distinct marks per tetrahedron. Every unmarked chamber lies between
two successive parallel normal discs of the same type and is a prism
$D^2\times I$. Its faces inherited from source faces are vertical rectangles;
its two horizontal faces lie on the cutting surface.

A corner region at a triangle/quad transition meets a quad-chain endpoint,
and a central face region also meets a marked endpoint. Thus gluing an
unmarked chamber only to other unmarked chambers matches vertical rectangles
between consecutive parallel arcs. The face identifications preserve or
reverse their interval direction. No global coorientation is assumed.

\begin{proposition}[Conservative cores and normal bundle midsections]
For a normal cut in the validated orientable source with $T$ tetrahedra,
at most $6T$ components touch the exceptional chambers. Every remaining
component is an interval bundle over a connected embedded normal midsection.
The marked components may themselves be interval bundles. The distinction
is sufficient rather than maximal.
\end{proposition}
\begin{proof}
Each marked component contains a distinct exceptional chamber, proving the
component bound. A component disjoint from the marks is assembled entirely
from the stated prisms, glued along their vertical rectangles. Their interval
fibres match, possibly with a reversal, giving an interval bundle. This is
the usual parallelity construction; no claim is made that our selected
components constitute the full maximal parallelity bundle.

Choose compatible middle points of the interval fibres in these prisms. In each tetrahedron they
form normal triangles or quadrilaterals of the original occupied type.
Their face arcs match because the vertical rectangle gluing identifies the
same two bounding normal arcs. Hence they form an embedded normal midsection.
Projecting the connected bundle to its base proves that this surface is
connected. Reversing fibre transitions can give a nonorientable midsection
and a twisted interval bundle; it does not obstruct the construction.
\end{proof}

\begin{corollary}[Unpatterned component-type bound]
Let $v$ be the number of source quotient vertices and $q$ the number of
positive quadrilateral types in the cutting vector. The cut has at most
$6T+v+2q\le12T$ distinct abstract unpatterned component homeomorphism types.
\end{corollary}
\begin{proof}
There are at most $6T$ marked components. The normal midsections of the
unmarked components are simultaneously disjoint, use only the original
$q$ quadrilateral types, and form one admissible normal surface by union.
The component-type theorem of Section~\ref{sec:compact-coordinates} therefore
bounds their distinct full normal vectors by $v+2q$. Equal vectors give
normally isotopic bases. Since the total manifold is orientable, its interval
bundle is the product over an orientable base and the orientable twisted
bundle over a nonorientable base. Thus the base determines the unpatterned
bundle type. Finally $v\le4T$ and $q\le T$.
\end{proof}

This bounds types, not total multiplicity, marked-component cell complexity
or the complexity of embeddings and attachment patterns. It does not
construct a compressed inventory or prove that arbitrary labelled
attachments have only linearly many types.

\subsection{Count marked and prismatic pieces without expansion}

Let $C$ be the original chamber-orbit count. Add singleton pairings joining
all exceptional chamber points to point zero, which is itself marked.
If $K$ original orbits touch marks, the new cone merges exactly those $K$
orbits and leaves every unmarked orbit separate. Its orbit count is
\[
                 C_\mathrm{cone}=C-K+1.
\]
Therefore the core count is $K=C-C_\mathrm{cone}+1$ and the wholly prismatic
count is $C_\mathrm{cone}-1$. The cone adds fewer than $6T$ pairings to the
same binary universe, so both queries retain polynomial bit cost.

The opt-in API flag is \code{classify\_prisms}. Default false preserves the
original scalar count and its version-one proof. Version two binds both
ordinary orbit traces and the conservative classification to the same
normalized source/vector fingerprint. The independent checker reconstructs
both the face-region system and the exceptional points; it imports neither
mark producer nor orbit planner. It verifies both traces, the cone identity
and the $6T$ bound. Cycle and replay-operation limits are shared across the
two queries. An unfinished classification has no top-level count or proof.

The result field \code{prismatic\_components} includes both product and
twisted interval bundles. It is not an assertion that all bases are
orientable, nor a maximal interval-bundle decomposition. No full cut
geometry, pattern transport, component extraction or knot verdict is added.

\subsection{Finite geometric and binary proof evidence}
All 1,555 maintained tests pass in 287.156 seconds; the 15 focused tests pass
in 0.598 seconds. The seven new regressions check scaled actual sources,
meridian/link/one-sided families, exact frozen disabled behavior, legacy
proof replay, shared cycle and operation caps, malformed cone/classification
data and cancellation. The final API uses ``prismatic'' rather than claiming
that every interval bundle is a Cartesian product.

The finite geometric pilot starts from the 1,635 earlier primitive or empty
cases and adds parallel multiples two, three and five where the primitive
disc count is at most 32. There are 5,241 cut-surface instances on 48 source
triangulations. An explicit small chamber graph locates its unmarked components;
each component's prism counts give a normal midsection satisfying the original
matching equations. All 7,838 such midsections are connected according to
Regina, and their interval-bundle Euler/boundary signatures occur with the
required multiplicities among the actual Regina cut pieces.
These signature checks are finite geometric consistency checks, not an
independent complete 3-manifold homeomorphism test. The interval-bundle
conclusion comes from the prism gluing proof above.

The integrated 61.696-second audit matches all 5,241 count classifications
and preserves the entire scalar output exactly against the frozen preceding
implementation when classification is disabled. Every classification proof
replays with mark, chamber and orbit producers disabled. All 157 preceding
version-one certificates also replay. Seven large meridian cases on one-,
four- and sixteen-tetrahedron sources use scales $2^{128}$, $2^{500}$ or
$2^{16384}$. Each has one marked core and $m-1$ prismatic pieces. The
one-tetrahedron family takes eight total orbit cycles, including the
16,385-bit case; the sixteen-tetrahedron family takes 136 cycles at both
large scales. The audit retains 55 representative version-two source proofs.

Five paired timing rounds measure 80 completed classification calls and
sixteen warmups, with A/A controls. Both arms construct the same valid
one-tetrahedron meridian source and compute total, marked and wholly prismatic
counts. The explicit reference expands the chamber graph, performs union-find
and tests the marked roots. It is an expansion baseline, not an optimized
Regina or published parallelity-bundle implementation. The native arm runs
both compressed orbit queries, independently replays its complete certificate
and serializes the result. All preparation and result serialization are
included. Neither arm extracts full core geometry or a component inventory.
\begin{center}
\begin{tabular}{rrrrrrr}
\toprule
copies & expanded points & explicit ms & native ms & ratio & old A/A & new A/A\\
\midrule
16 & 49 & 0.148 & 0.632 & 0.230 & 0.976 & 1.015\\
256 & 769 & 0.410 & 0.660 & 0.611 & 0.976 & 0.995\\
4096 & 12289 & 4.720 & 0.677 & 7.062 & 1.013 & 0.980\\
65536 & 196609 & 83.157 & 0.748 & 112.287 & 1.018 & 0.907\\
\bottomrule
\end{tabular}
\end{center}


At 65,536 meridians the explicit graph has 196,609 points. Its median
classification time is 83.157 milliseconds, versus 0.748 milliseconds for
the native method; the paired ratio is 112.287. At 4,096 copies the ratio
is 7.062. At sixteen and 256 copies the ratios are 0.230 and 0.611:
proof production and independent native replay lose on these small cases.
The old A/A ratios range from 0.976 to 1.018 and the new ones from 0.907
to 1.015. These are host observations, not precise universal speedups.
The asymptotic benefit is avoiding a graph with one node per represented
chamber while returning the same conservative classification.

The default scalar operation, the complete recognizer and the optional
geometric recognition defaults remain unchanged. No full-recognition speedup
or general quasi-polynomial bound follows from these count-task timings.
The existence of linearly many abstract unpatterned types does not extract
or authenticate their patterned representations.

Runtime release \code{eced40c56}, the complete test logs, finite geometric
pilot, source-proof audit and timing records are retained.

The driver is \code{normal\_orbit\_research.cut\_products}. Records are
under \code{data/cut-products-*} and \code{data/cut-product-pilot.*}.
Publication review checks 660 source/evidence hashes, the exact pilot inputs,
all timing outcomes and the passing test records. It independently replays
all 55 saved version-two source proofs with producers disabled.
The next step is a source-authenticated compressed inventory and core geometry,
with boundary patterns and attachment data retained; those obligations remain
open rather than being replaced by a count or a type bound.

